We investigate article-level political ideology classification in U.S. news media through three complementary approaches: metalinguistic feature integration, large language model (LLM) evaluation and fine-tuning, and causal inference analysis. We introduce a multi-objective continued pre-training approach that incorporates article-level theme and tone prediction as auxiliary tasks alongside contrastive triplet-loss objectives. After fine-tuning on the AllSidesV2 dataset, our best domain-adapted RoBERTa configuration achieves 62.19 F1-Macro, outperforming the previous state-of-the-art POLITICS model and maintaining superior performance over zero-shot LLMs. Fine-tuned GPT-4o-mini surpasses our transformer model with sufficient training data (69.37 F1-Macro with 2,000 samples), though at the cost of potential source leakage from its opaque pre-training corpus.

We then apply a causal inference framework using Double Machine Learning and mediation analysis to examine whether topic sentiment causally affects perceived political ideology, comparing five annotation paradigms: expert human labels, fine-tuned RoBERTa, zero-shot GPT-4o-mini, fine-tuned GPT-4o-mini, and Llama-3.3-70B. At the community level, only the fine-tuned GPT-4o-mini, despite achieving the highest classification accuracy (F1 = 72.48), produces significant causal effects for Social Issues, Geopolitics, and Trump Presidency communities, while human and Llama-3.3-70B annotations remain null throughout. Llama-3.3-70B exhibits a distinct failure mode of systematic effect attenuation toward zero. Mediation analysis reveals that the fine-tuned GPT effect operates almost entirely through a direct pathway (96.8\% of the total effect), consistent with shortcut learning: fine-tuning introduces spurious sentiment-ideology coupling that is invisible to standard F1-based evaluation but surfaces under causal probing. At the topic level, no robust causal effects emerge, and estimated causal directions can reverse depending on the annotation model. These findings demonstrate mediation analysis as a diagnostic framework for annotation trustworthiness and caution against using LLM annotations as drop-in replacements for human labels in downstream causal analyses.

We also release AllSidesV2, an expanded version of the AllSides dataset with over 10,000 additional annotated articles to facilitate future research in political ideology detection.

